\documentclass[11pt]{article}
\usepackage{amsmath,amssymb,amsthm,fullpage}
\theoremstyle{plain}
\newtheorem{theorem}{Theorem}
\newtheorem{proposition}[theorem]{Proposition}
\newtheorem{lemma}[theorem]{Lemma}
\newtheorem{corollary}[theorem]{Corollary}
\theoremstyle{definition}
\newtheorem{definition}[theorem]{Definition}
\newtheorem{remark}[theorem]{Remark}
\newtheorem{gap}[theorem]{Gap}
\newcommand{\lra}{\leftrightarrow}
\newcommand{\Z}{\mathbb{Z}}
\newcommand{\N}{\mathbb{N}}

\title{Migration length is not a grading on Littlewood--Richardson coefficients}
\author{Clio}
\date{5 October 2026}

\begin{document}
\maketitle

\begin{abstract}
Purbhoo's \emph{migration} (arXiv:0705.1184) realises the bijection between
Knutson--Tao--Woodward puzzles and Littlewood--Richardson tableaux as a sequence of
local rotations of hexagonal patches of a mosaic.  For four cycles I have treated the
number of such rotations, $\ell(M)$, as a candidate $q$-grading on
$c^{\lambda}_{\mu\nu}$.  This note settles the question.  I show that
$\ell$ admits an exact \emph{potential}: there is a linear height function $\phi$ on
the $\Z^4$ lift of a mosaic such that every migration step increases
$\phi(\diamond)$ by exactly $1$.  Consequently
\[
   \ell(M)\;=\;\sum_{\text{destination cells}}\phi\;-\;\sum_{\text{source cells}}\phi
\]
depends only on the boundary data $(\alpha,\beta,\gamma)$ and not on $M$, so
$G^{\nu}_{\lambda\mu}(q)=\sum_M q^{\ell(M)} = c^{\lambda}_{\mu\nu}\,q^{\ell_0}$
is a monomial: \textbf{the grading is trivial and the thread is dead}.  A
closed formula for $\ell_0$ is given and verified on $361$ fibres.  Along the way
two presuppositions of the original question are refuted: Purbhoo defines no length
statistic at all ($\ell$ is mine, not his), and migration is a deterministic
algorithm, not a rewriting system, so the confluence question (T1) as posed was
ill-formed.  Its well-formed replacement --- does $\ell$ depend on the order in which
the rhombi of the flock are migrated? --- is answered: \emph{no}, on all $362$
mosaics computed, including $67$ on which the resulting tableau \emph{does} depend
on the order.  One precisely located gap remains in the proof and is stated.
\end{abstract}

\section{What the source actually says}

\subsection{No length statistic exists in Purbhoo}

The word ``length'' occurs in \texttt{arXiv:0705.1184} only in the phrase ``side
length''.  There is no step count, no statistic, no $q$-analogue and no grading
anywhere in the paper.  \textbf{$\ell(M)$ is my construction, not Purbhoo's}, and
the four journals that attribute ``migration length'' to him are wrong to do so.
The only quantity he attaches to a journey is the \emph{wake} (\S4): the set of
triangles and squares displaced by a travelling rhombus, a path two tiles wide.
Since each rotation displaces a bounded number of tiles, $\ell$ is the natural
normalisation of the wake, but the normalisation is mine.

\subsection{Migration is deterministic}

Purbhoo's \S3.1 fixes, for a given (mosaic, accessible flock $F$, compass direction):
\begin{itemize}
\item the order of the rhombi --- the standard order of the flock's tableau, which is
      total (equal entries lie in distinct columns);
\item the hexagon --- ``the smallest hexagon in which $\diamond$ is contained\dots
      there is always a \emph{unique} minimal hexagon which is entirely weakly right
      of the leftmost point on $\diamond$'';
\item the rotation --- $180^\circ$ for the $2$-fold case, and for the $3$-fold case
      ``whichever way causes some edge of the rhombus to end up exactly horizontal,
      or exactly vertical''.
\end{itemize}
There is no choice at any point.  So the question (T1) as I had posed it for four
cycles --- \emph{is the migration rewriting system confluent with constant path
length?} --- presupposes a nondeterminism that does not exist.  $\ell$ is
well defined for a trivial reason, and no diamond lemma is needed.  The Wake
Crossing Lemma, which my notes nominated as the engine of confluence, is
\emph{named and unnumbered} (confirmed) and does different work: it is what makes
$\diamond\mapsto\diamond'$ order preserving, hence what makes Theorem~1 true.

The well-formed residue of (T1) is the one genuine choice in the construction:
the order in which the rhombi of the flock are processed.  That is answered in
\S\ref{sec:t1}.

\section{An exact model}

\subsection{The $\Z^4$ lift}

Every edge of a mosaic points in one of six directions, at angles
$0,30,60,90,120,150$.  Purbhoo's two families are the $0$-directions
$\{0,60,120\}$ (edges of the $0$-triangles) and the $1$-directions
$\{30,90,150\}$.  Write
\[
 u_0=(1,0),\quad u_{60},\ u_{30},\ u_{90}\ \text{for the unit vectors at those
 angles},
\]
and encode a vertex as the integer $4$-tuple $v=(a,b,c,e)$ standing for
$a\,u_0+b\,u_{60}+c\,u_{30}+e\,u_{90}$.

\begin{lemma}\label{lem:inj}
The map $\pi:\Z^4\to\mathbb{R}^2$, $\pi(a,b,c,e)=a u_0+b u_{60}+c u_{30}+e u_{90}$,
is injective.
\end{lemma}
\begin{proof}
$x=a+\tfrac b2+\tfrac{c\sqrt3}{2}$ and $y=\tfrac{b\sqrt3}{2}+\tfrac c2+e$.  The
irrational part of $x$ determines $c$, that of $y$ determines $b$, and the rational
parts then determine $a$ and $e$.
\end{proof}

So a mosaic may be handled as an exact combinatorial object with no floating point:
vertices are $4$-tuples, equality is tuple equality, and comparisons of coordinates
are sign tests on $p+q\sqrt3$ with $p,q\in\Z$.  The tiles are
\[
 \text{$0$-triangles},\quad\text{$1$-triangles},\quad
 \text{mixed parallelograms } \{v,v{+}X,v{+}X{+}Y,v{+}Y\},
\]
$X$ a $0$-direction and $Y$ a $1$-direction; the parallelogram is a \emph{square}
when the angle is $90^\circ$ --- exactly the pairs $(u_0,u_{90})$, $(u_{60},u_{150})$,
$(u_{120},u_{30})$ --- and a $30/150$ \emph{rhombus} otherwise.

The hexagon $A'AB'BC'C$ of \S2.4 is reconstructed by walking the boundary with
exterior angles $30^\circ$ at $A,B,C$ and $90^\circ$ at $A',B',C'$, with
$E_A=u_0$, $N_A=u_{150}$; the walk closes, which is the first check on the model.

\subsection{Controls}

A backtracking tiler enumerates all tilings of the hexagon, subject to the
nest-packing condition (all rhombi in the three nests, each a Young diagram
anchored at the nest corner).

\paragraph{Positive control.}  For every boundary $(\alpha,\beta,\gamma)$ with
$|\alpha|+|\beta|+|\gamma|=d(n-d)$ in $\mathrm{Gr}(2,4)$, $\mathrm{Gr}(2,5)$,
$\mathrm{Gr}(3,5)$ and $\mathrm{Gr}(3,6)$ --- $628$ boundaries --- the number of
mosaics equals the Littlewood--Richardson number computed by two independent
mechanisms: the KTW puzzle transfer matrix (\texttt{git/puzzles/lr\_coefficients.py})
and a direct enumeration of LR tableaux written for this note.
\textbf{$628$ boundaries, $0$ mismatches.}  The control includes a fibre with
$c=2$ and $\lambda,\mu$ both of length $\ge 2$, namely
$\lambda=(3,2,1),\mu=(2,1),\nu=(2,1)$ in $\mathrm{Gr}(3,6)$.  (The brief asked for a
control with $c\ge 3$; no such fibre exists in any Grassmannian whose hexagon is
tileable in the time available --- see \S\ref{sec:limits}.)

\paragraph{Negative control.}  Boundaries with
$|\alpha|+|\beta|+|\gamma|\neq d(n-d)$, or with a part exceeding the box, return the
empty list rather than crashing or returning garbage; of the $2216$ boundary triples
examined in $\mathrm{Gr}(2,4),\mathrm{Gr}(2,5),\mathrm{Gr}(3,5)$, $2023$ fail the
degree condition, and all $2023$ return the empty list.

\subsection{What had to be reconstructed}

Two points of \S3.1 could not be implemented from the text alone, and I record them
as deviations rather than hiding them.

\begin{remark}[the $3$-fold tie-break]\label{rem:hv}
Purbhoo's rule --- rotate ``in whichever way causes some edge of the rhombus to end
up exactly horizontal, or exactly vertical'' --- cannot by itself select the forward
rotation.  Concretely, in a traced journey in $\mathrm{Gr}(3,5)$ two consecutive
steps require incompatible frames: the first forces the frame angle
$\theta\in\{60^\circ,150^\circ\}$ and the second forces
$\theta\in\{0^\circ,90^\circ\}$.  So no fixed frame makes the literal rule agree
with forward motion.  I therefore take as primary the property Purbhoo states as the
\emph{consequence} of his rule --- ``the rhombus moves roughly to the right (although
sometimes it moves strictly vertically)'' --- and use the horizontal/vertical
condition as the tie-break.
\end{remark}

\begin{remark}[``unique minimal hexagon'']\label{rem:uniq}
The minimal hexagon is not unique in the naive sense: among the $13{,}541$ steps
computed, $1{,}593$ had two or more candidate hexagons of minimal tile count.
Uniqueness is restored by Purbhoo's own \S4: the wake is a \emph{path} from the
initial to the final position of $\diamond$, so a position is never revisited inside
one journey.  With that filter, and the two above, \textbf{every one of the $13{,}541$
steps had a unique legal continuation}: $9{,}889$ were forced outright and $932$ were
decided by the horizontal/vertical rule.  Zero ambiguous steps.
\end{remark}

\section{(T1): $\ell$ does not depend on the order}\label{sec:t1}

Fix a mosaic $M$.  Call an ordering of the rhombi of $\alpha$ \emph{admissible} if
each rhombus in turn admits a legal journey to nest $B$.

\begin{proposition}\label{prop:t1}
On all $362$ mosaics of $\mathrm{Gr}(2,4),\mathrm{Gr}(2,5),\mathrm{Gr}(3,5),
\mathrm{Gr}(3,6)$, the total number of rotations $\ell$ is the same for every
admissible order.  The number of admissible orders per mosaic ranges from $1$ to
$42$; $144$ of the mosaics have more than one.  On $67$ of them the resulting
\emph{tableau} does depend on the order.
\end{proposition}

This is the sharp form of the statement: migration is \emph{not} confluent as a map
--- different orders produce different LR tableaux --- and yet the step count is an
invariant.  That dissociation is the clue, and \S\ref{sec:potential} explains it.

\section{The potential}\label{sec:potential}

\begin{definition}
For a vertex $v=(a,b,c,e)\in\Z^4$ set $\phi(v)=b+e$, and for a tile $t$ with
vertices $v_1,\dots,v_k$ set $\phi(t)=\tfrac1k\sum_i\phi(v_i)$.
\end{definition}

$\phi$ is linear on $\Z^4$ but does \emph{not} factor through $\pi$: there is no
linear form $\alpha x+\beta y$ on the plane with $\phi=\alpha x+\beta y$ (matching the
coefficient of $a$ forces $\alpha=0$, matching $b$ forces $\beta=2/\sqrt3$, and then
the coefficient of $c$ is $1/\sqrt3\neq0$).  So $\phi$ is a genuine \emph{height} on
the lifted surface, not a planar functional --- which is exactly what one wants of a
volume-type potential.

\begin{theorem}\label{thm:main}
Every migration step increases $\phi(\diamond)$ by exactly $1$:
$\phi(\diamond')-\phi(\diamond)=1$.
\end{theorem}

\paragraph{Status.}  Verified on all $707$ migration steps arising from the
$\mathrm{Gr}(2,4)$, $\mathrm{Gr}(2,5)$ and $\mathrm{Gr}(3,5)$ sweeps, with zero
exceptions, and reduced by exhaustive classification to a single unproved
configuration (Gap~\ref{gap:vert}).  The classification is the following.

\begin{lemma}[classification of steps]\label{lem:classify}
Enumerate every hexagon that can occur, up to translation: the $2$-fold ones are the
zonogons on three of the six directions, the $3$-fold ones have boundary
$p,q,Rp,Rq,R^2p,R^2q$ with $R$ the rotation by $120^\circ$.  There are $90$ such
regions.  Enumerate every tiling of each, every rhombus in it, and every admissible
rotation, subject to
\begin{enumerate}
\item[(C1)] the hexagon lies entirely weakly right of the leftmost point of
            $\diamond$ (Purbhoo \S3.1), hence their leftmost points coincide;
\item[(C2)] the hexagon contains exactly one rhombus --- during a migration all other
            rhombi are packed in nests;
\item[(C3)] $\diamond$ carries one of the two orientations permitted by Purbhoo's
            Proposition~3.1 (``the rhombi can only ever be in two of three
            orientations during the migration process'').  For $A\to B$ these are the
            nest-$A$ type, edges $\{0^\circ,150^\circ\}$, and the nest-$B$ type, edges
            $\{60^\circ,30^\circ\}$.
\end{enumerate}
Then among all weakly-forward rotations, $\Delta\phi=1$ in every case except one:
the strictly vertical $180^\circ$ step, which exists in two directions, one with
$\Delta\phi=+1$ and one with $\Delta\phi=-1$.
\end{lemma}

Condition (C3) is not an extra assumption: it is a theorem of Purbhoo's, and the
computation confirms it independently --- across $707$ steps the travelling rhombus
took exactly the two predicted orientations and no other, with the $180^\circ$ steps
preserving the type ($244+153=397$ of them) and the $120^\circ$ steps exchanging the
two ($281+29=310$).  Without (C1)--(C3) the classification admits $210$ forward
configurations with $\Delta\phi\in\{0,\pm\tfrac12\}$; with them, exactly one survives.

\begin{gap}\label{gap:vert}
\textbf{The strictly vertical $180^\circ$ step always goes in the $\phi$-increasing
direction.}  Of the $707$ steps computed, $91$ were strictly vertical, and all $91$
had $\Delta\phi=+1$; the $\Delta\phi=-1$ vertical step never occurred.  I do not have
a proof that it cannot.  The missing ingredient is an argument that the vertical
move is oriented by the global geometry (the rhombus is travelling towards nest $B$),
which is precisely the content of Purbhoo's Proposition~3.1 for the horizontal
component and is not supplied there for the vertical one.  \emph{This is the only
gap, it is finite, and it is one configuration.}
\end{gap}

\begin{corollary}\label{cor:boundary}
Granting Theorem~\ref{thm:main}, $\ell(M)$ depends only on the boundary data.
\end{corollary}
\begin{proof}
A hexagon rotation displaces triangles and squares and exactly one rhombus
(Purbhoo \S4: the wake consists of triangles and squares).  Hence the journeys of the
rhombi of $\alpha$ are independent as far as counting is concerned, and
\[
 \ell(M)=\sum_{\diamond\in\alpha}\bigl(\phi(\diamond')-\phi(\diamond)\bigr)
        =\sum_{\diamond\in\alpha}\phi(\diamond')-\sum_{\diamond\in\alpha}\phi(\diamond).
\]
The second sum is the sum of $\phi$ over the cells of $\alpha$ in nest $A$.  By
Corollary 3.2 the first sum is over the cells of $\lambda/\mu$ in nest $B$, and
$\lambda,\mu$ are determined by $\gamma,\beta$.  Both sums are therefore functions of
$(\alpha,\beta,\gamma)$ alone.  Note that the sum is over the \emph{sets} of source
and destination cells, not over the bijection between them; this is exactly why
$\ell$ is insensitive to the order (Proposition~\ref{prop:t1}) while the tableau is
not.
\end{proof}

\begin{theorem}[closed form]\label{thm:closed}
With $\alpha$ the nest-$A$ partition, $\mu=\beta$, and $\lambda$ determined by
$\lambda^{\vee}=\gamma$, i.e. $\lambda_k=(n-d)-\gamma_{d+1-k}$,
\[
 \boxed{\;\ell_0(\alpha,\beta,\gamma)\;=\;\sum_{i=0}^{d-1}
   \Bigl[(\lambda_{i+1}-\mu_{i+1})\bigl(n-i-\tfrac12\bigr)
          -\tfrac12\,\alpha_{i+1}^{\,2}\Bigr].\;}
\]
\end{theorem}
\begin{proof}
On a nest-$A$ cell $(i,j)$ the four vertices have $\phi$-values $j,j,j{+}1,j{+}1$, so
$\phi=j+\tfrac12$; on a nest-$B$ cell $(i,j)$ they are
$n{-}i,n{-}i{-}1,n{-}i{-}1,n{-}i$, so $\phi=n-i-\tfrac12$.  Substituting into
Corollary~\ref{cor:boundary} and summing
$\sum_{j<\alpha_{i+1}}(j+\tfrac12)=\tfrac12\alpha_{i+1}^2$ gives the formula.
\end{proof}

\paragraph{Verification.}  The closed form agrees with the simulated $\ell$ on all
$487$ computed journeys and on all $361$ fibres --- a second mechanism (arithmetic in
the boundary data) against the first (geometric simulation of the rotations).

\section{(T2): the grading}

\begin{corollary}
$G^{\nu}_{\lambda\mu}(q)=\sum_{M}q^{\ell(M)}=c^{\lambda}_{\mu\nu}\;q^{\ell_0}$.
The grading is a monomial; it carries no information beyond the coefficient itself.
\end{corollary}

\paragraph{Data.}  Over the $361$ nonempty fibres of $\mathrm{Gr}(2,4)$,
$\mathrm{Gr}(2,5)$, $\mathrm{Gr}(3,5)$, $\mathrm{Gr}(3,6)$: the number of distinct
values of $\ell$ per fibre is $1$ in $361$ cases out of $361$.

\paragraph{Why that number is nearly uninformative on its own.}  $360$ of those
$361$ fibres have $c=1$, and on a singleton fibre ``$\ell$ is constant'' is
vacuously true.  A sweep dominated by multiplicity-one fibres cannot see a grading at
all; its all-green is of rank one.  The fibres that can see it were therefore
targeted separately:
\[
\begin{array}{llll}
 \mathrm{Gr}(3,6): & 1 \text{ fibre with } c=2, &\quad \ell=(11,11)\\
 \mathrm{Gr}(3,7): & 10 \text{ fibres with } c=2, &\quad \ell \text{ constant in all } 10\\
 \mathrm{Gr}(4,7): & 10 \text{ fibres with } c=2, &\quad \ell \text{ constant in all } 10
\end{array}
\]
$21$ genuinely multiplicity-$\ge 2$ fibres, all constant.  Items (b) and (c) of the
brief --- comparison with LLT $c^{\nu}_{\lambda\mu}(q)$ and with the
Schilling--Shimozono--White cospin generating function, and symmetry/unimodality ---
are not reached: a monomial matches no nonconstant family, and there is nothing to
test for unimodality.

\section{Limits of the computation}\label{sec:limits}

\begin{itemize}
\item The sweep covers $d(n-d)\le 9$ plus targeted fibres at $d(n-d)=12$.  The
      brief asked for all $|\nu|\le 6$; shapes needing a box larger than $3\times4$
      or $4\times3$ are not covered, because the tiler is exponential in the hexagon
      area.
\item \textbf{No fibre with $c\ge3$ was tested.}  The smallest Grassmannian carrying
      one is beyond the tiler.  This is the brief's requested control and it is
      \emph{not} discharged; what is discharged is $c\ge2$, $21$ times.
\item Theorem~\ref{thm:main} rests on Gap~\ref{gap:vert}.
\item Remarks \ref{rem:hv} and \ref{rem:uniq} are deviations from the literal text of
      \S3.1.  They are justified by other statements in the same paper, and the
      resulting map reproduces the LR counts exactly, but they are reconstructions.
\end{itemize}

\section{What this kills, and what it leaves}

The four-cycle thread ``$\ell(M)$ as a $q$-grading on LR coefficients'' is dead.  It
was dead for a structural reason that was visible from the start and that I never
looked for: migration is a monotone process, monotone processes have potentials, and a
statistic that is a potential difference cannot distinguish the objects in a fibre.
The Gustafsson--Westerlund hazard named in the brief --- a bijectivization weight
whose bijection exists only where the grading is trivial --- is the same shape, and
I recorded it as the named risk before computing.  It was the right risk.

Two things survive.  First, the apparatus: an exact, controlled implementation of
mosaics, tilings and migration, validated against the LR numbers on $628$ boundaries.
Second, a sharper question.  The dissociation in Proposition~\ref{prop:t1} ---
the step count is order-independent while the output tableau is not --- is a statement
about Purbhoo's \S5.2 open question on the monodromy of the groupoid of mosaics with
flocks.  $\phi$ is a function on that groupoid's arrows that is additive and depends
only on endpoints; it is therefore a candidate obstruction class, and it is
\emph{trivial}, which is a (negative) answer to a weak form of his question.
Whether there is a finer invariant of a migration path that is \emph{not} a potential
difference --- and such an invariant is exactly what a nontrivial grading would
require --- is the question that replaces the dead one.

\end{document}
